\expandafter\gdef\csname SNside1\endcsname{\node[anchor=north west,inner sep=0pt,text width=48mm,font=\SideFont,align=left] at ([xshift=138mm,yshift=-516.15pt]current page.north west) {\textcolor{SNink}{\textbf{错项排除 · 角色由这次交互决定}}\par\vspace{1.5mm}一台机器可向DNS发查询，同时给其他客户提供Web服务。它在前一次交互中是客户，在后一次交互中是服务器；设备性能、是否长期在线以及机箱外观都不能替代应用角色的判断。};}
\expandafter\gdef\csname SNside2\endcsname{\node[anchor=north west,inner sep=0pt,text width=48mm,font=\SideFont,align=left] at ([xshift=138mm,yshift=-478.22pt]current page.north west) {\textcolor{SNink}{\textbf{课程连接 · 分发下界就是资源守恒}}\par\vspace{1.5mm}操作系统调度题用总工作量除可用服务率找下界，文件分发也一样。服务器上传、最慢客户下载和系统总上传分别给出限制，取最大值才能同时满足；它们不是三个必须顺序相加的阶段。};}
\expandafter\gdef\csname SNside3\endcsname{\node[anchor=north west,inner sep=0pt,text width=48mm,font=\SideFont,align=left] at ([xshift=138mm,yshift=-95.00pt]current page.north west) {\textcolor{SNink}{\textbf{最小反例 · 一个域可以跨多个区}}\par\vspace{1.5mm}abc.com把子域lab.abc.com委派给另一组权威服务器后，lab仍在abc.com名称子树内，却已不属于父区的直接权威数据范围。这个委派例子足以排除``域和区永远完全相同''的说法。};}
\expandafter\gdef\csname SNside4\endcsname{\node[anchor=north west,inner sep=0pt,text width=48mm,font=\SideFont,align=left] at ([xshift=138mm,yshift=-95.00pt]current page.north west) {\textcolor{SNink}{\textbf{不变量 · 每条查询都有发出者}}\par\vspace{1.5mm}主机向本地解析器的一条请求，不是本地解析器发出的请求。题目问某台服务器发出多少次时，应逐条画箭头并只数它发出的边；缓存命中可使远程次数归零，却不改变递归与迭代的定义。};}
\expandafter\gdef\csname SNside5\endcsname{\node[anchor=north west,inner sep=0pt,text width=48mm,font=\SideFont,align=left] at ([xshift=138mm,yshift=-95.00pt]current page.north west) {\textcolor{SNink}{\textbf{最小反例 · 数据端口不总是20}}\par\vspace{1.5mm}FTP被动模式由服务器通告一个监听数据端口，再由客户连接该端口。控制仍走服务器21端口，数据却可走另一个值，因此``服务器所有数据连接端口恒为20''已被这个合法模式反驳。};}
\expandafter\gdef\csname SNside6\endcsname{\node[anchor=north west,inner sep=0pt,text width=48mm,font=\SideFont,align=left] at ([xshift=138mm,yshift=-95.00pt]current page.north west) {\textcolor{SNink}{\textbf{错项排除 · 编码不提供保密性}}\par\vspace{1.5mm}Base64把任意字节表示成ASCII字符，接收者可按公开规则还原。它解决邮件内容表示与传送兼容，不保证只有特定收件人可读；MIME也不替代SMTP投递或POP3读取机制。};}
\expandafter\gdef\csname SNside7\endcsname{\node[anchor=north west,inner sep=0pt,text width=48mm,font=\SideFont,align=left] at ([xshift=138mm,yshift=-95.00pt]current page.north west) {\textcolor{SNink}{\textbf{最小反例 · 持久连接仍可逐个等待}}\par\vspace{1.5mm}已建立一条TCP连接，先请求图片A，等它完整返回后才请求B，这是持久但非流水。两对象仍付两次请求响应等待，所以``用了持久连接就所有对象只要1 RTT''的说法不成立。\par\smallskip\begin{tikzpicture}[x=1mm,y=1mm,>=Stealth,node distance=5mm]
\node[draw=SNink!50,fill=SNblue,rounded corners=1mm,text width=34mm,align=center,inner sep=1.2mm,font=\fontsize{7.5}{9}\selectfont] (m0) at (0,0) {同一TCP};
\node[draw=SNink!50,fill=SNblue,rounded corners=1mm,text width=34mm,align=center,inner sep=1.2mm,font=\fontsize{7.5}{9}\selectfont] (m1) at (0,-11) {请求A};
\node[draw=SNink!50,fill=SNblue,rounded corners=1mm,text width=34mm,align=center,inner sep=1.2mm,font=\fontsize{7.5}{9}\selectfont] (m2) at (0,-22) {等A返回};
\node[draw=SNink!50,fill=SNblue,rounded corners=1mm,text width=34mm,align=center,inner sep=1.2mm,font=\fontsize{7.5}{9}\selectfont] (m3) at (0,-33) {再请求B};
\draw[->,SNink] (m0) -- (m1);
\draw[->,SNink] (m1) -- (m2);
\draw[->,SNink] (m2) -- (m3);
\end{tikzpicture}
\par{\fontsize{7.5}{9}\selectfont 持久连接仍可以不使用流水}};}
\expandafter\gdef\csname SNside8\endcsname{\node[anchor=north west,inner sep=0pt,text width=48mm,font=\SideFont,align=left] at ([xshift=138mm,yshift=-267.50pt]current page.north west) {\textcolor{SNink}{\textbf{不变量 · 资源身份不等于一次往返}}\par\vspace{1.5mm}一份HTML可以引用多个图片和样式文件，且服务器可能不同。页面在界面中只有一个标签，不代表网络里只有一个对象；必须先发现资源地址，才有依据安排后续DNS、连接和请求。};}
\expandafter\gdef\csname SNside9\endcsname{\node[anchor=north west,inner sep=0pt,text width=48mm,font=\SideFont,align=left] at ([xshift=138mm,yshift=-95.00pt]current page.north west) {\textcolor{SNink}{\textbf{错项排除 · 流水不能猜未知资源}}\par\vspace{1.5mm}在HTML未返回且没有其他先验信息时，浏览器还不知道其中的图片URL。允许流水只能把发现后的多个请求重叠，不能把HTML阶段也凭空抹去；这正是理想模型仍保留三个RTT阶段的原因。};}
\expandafter\gdef\csname SNside10\endcsname{\node[anchor=north west,inner sep=0pt,text width=48mm,font=\SideFont,align=left] at ([xshift=138mm,yshift=-95.00pt]current page.north west) {\textcolor{SNink}{\textbf{不变量 · 应用数据量与FIN分开计}}\par\vspace{1.5mm}FTP文件18000 B，从SYN序号100开始，则首数据101，末数据18100，FIN为18101，对方确认FIN为18102。控制连接已有自己的序号空间，不能混进这条数据连接的累计字节数。};}
\expandafter\gdef\csname SNside11\endcsname{\node[anchor=north west,inner sep=0pt,text width=48mm,font=\SideFont,align=left] at ([xshift=138mm,yshift=-297.09pt]current page.north west) {\textcolor{SNink}{\textbf{最小反例 · 缓存命中会减少ARP广播}}\par\vspace{1.5mm}若本次已缓存Web服务器IP，就可直接解析网关MAC，不必先访问本地DNS。即使ARP表和交换表为空，也不足以证明必出现两次ARP广播；该结论还需要一次新的DNS查询这一条件。};}
\expandafter\gdef\csname SNside12\endcsname{\node[anchor=north west,inner sep=0pt,text width=48mm,font=\SideFont,align=left] at ([xshift=138mm,yshift=-95.00pt]current page.north west) {\textcolor{SNink}{\textbf{不变量 · 304仍是一份网络响应}}\par\vspace{1.5mm}条件请求发到服务器后收到304，说明发生了验证往返，只免去了完整实体内容。新鲜缓存直接命中则可能完全不发该请求。两种缓存行为都节省流量，却不具有相同的时延表达式。};}
\expandafter\gdef\csname SNside13\endcsname{\node[anchor=north west,inner sep=0pt,text width=48mm,font=\SideFont,align=left] at ([xshift=138mm,yshift=-95.00pt]current page.north west) {\textcolor{SNink}{\textbf{不变量 · 应用数据量与FIN分开计}}\par\vspace{1.5mm}FTP文件18000 B，从SYN序号100开始，则首数据101，末数据18100，FIN为18101，对方确认FIN为18102。控制连接已有自己的序号空间，不能混进这条数据连接的累计字节数。};}
\expandafter\gdef\csname SNside14\endcsname{\node[anchor=north west,inner sep=0pt,text width=48mm,font=\SideFont,align=left] at ([xshift=138mm,yshift=-441.86pt]current page.north west) {\textcolor{SNink}{\textbf{最小反例 · 缓存命中会减少ARP广播}}\par\vspace{1.5mm}若本次已缓存Web服务器IP，就可直接解析网关MAC，不必先访问本地DNS。即使ARP表和交换表为空，也不足以证明必出现两次ARP广播；该结论还需要一次新的DNS查询这一条件。};}
\expandafter\gdef\csname SNside15\endcsname{\node[anchor=north west,inner sep=0pt,text width=48mm,font=\SideFont,align=left] at ([xshift=138mm,yshift=-95.00pt]current page.north west) {\textcolor{SNink}{\textbf{不变量 · 读与发是不同应用动作}}\par\vspace{1.5mm}SMTP负责提交和转送，POP3负责读取邮箱。同一用户的一封邮件可能依次经过两种协议，不能因二者都处理邮件就互换。POP3会话能读取多封，下载后是否删除还取决于相应操作和更新阶段。};}
\AddToShipoutPictureFG{\begin{tikzpicture}[remember picture,overlay]\ifcsname SNside\thepage\endcsname\csname SNside\thepage\endcsname\fi\end{tikzpicture}}